\documentclass[12pt]{article}
\usepackage{graphicx}
\usepackage{amsmath}
\usepackage{amssymb}
\usepackage{amsthm}
\usepackage{url}
\usepackage{booktabs}
\usepackage[margin=1in]{geometry}
\usepackage{natbib}

\newtheorem{proposition}{Proposition}
\newtheorem{assumption}{Assumption}

\title{Cross-Border Demand, License Caps, and the Timing of Entry\\
in U.S.\ Cannabis Retail Markets}
\author{Mahyar Rezazadeh \\ Babak Sahragard}
\date{August 8, 2026}

\begin{document}
\maketitle

\begin{abstract}
\noindent We document three facts about cannabis retailers in U.S.\ counties located near the border of a
state where recreational cannabis remains illegal. First, revenue per dispensary is significantly higher in
these border counties, both at a 70-mile threshold and continuously in distance to the nearest non-legal
county. Second, this pattern is consistent with a license-cap mechanism: in Colorado and Washington, where
the number of retail licenses is constrained by state or local policy, the border effect appears in revenue
per store rather than in the number of stores or in revenue per capita; in Montana, the one state in our
sample with no statewide numerical cap, the same border proximity instead raises revenue per capita while
leaving revenue per store unaffected. Third, using each dispensary's own opening date across a nine-state
sample, we show that a county's \emph{second} dispensary opens significantly faster the closer that county
is to a non-legal border, even though its first dispensary does not, and that this effect fades by the third
and fourth entrant. We then develop a simple entry game, in the spirit of Bayesian learning models of
industry evolution, in which a first entrant's regulatory delay is unrelated to local demand, but its
operation reveals a signal about previously uncertain cross-border demand that a potential second entrant
conditions on. The model reproduces all three timing facts and explains, rather than merely reports, why a
classical cross-sectional entry-threshold test does not show the same distance gradient the timing result
does.
\end{abstract}

\section{Introduction}

A recurring question in the economics of regulated retail markets is how demand from outside a market's
official boundaries -- a nearby jurisdiction with different tax or legal treatment of the same good --
affects entry, pricing, and market structure inside that market. We study this question using U.S.\
recreational cannabis retail, where state-level legalization has proceeded unevenly and many legal-market
counties sit directly adjacent to a state where cannabis remains illegal. \citet{HansenMillerWeber2020} and
\citet{HaoCowan2020} establish that cross-border cannabis demand is real: Washington retailers near the
Oregon border lost a substantial share of sales once Oregon's own market opened, and marijuana-related
arrests respond to legal availability across a border. We build on this evidence in three steps and then
develop a simple theoretical account of the third.

\emph{First}, we show that revenue per dispensary is higher in counties near a non-legal border, both using
a 70-mile threshold and a continuous distance measure (Section~\ref{sec:result1}). \emph{Second}, we ask what
mechanism can explain this and argue, using Montana as a within-sample contrast, that state licensing caps
are doing the work: where entry is capped, captive demand shows up as revenue per existing store; where it
is not, it shows up as more stores instead (Section~\ref{sec:result2}). \emph{Third}, we turn to entry
\emph{timing} rather than revenue levels: using each dispensary's own opening date, we show a county's second
dispensary arrives significantly faster the closer that county is to a non-legal border, even though its
first dispensary does not, and that the effect is gone by the third and fourth entrant
(Section~\ref{sec:result3}). \emph{Fourth}, we build a simple entry game that explains this specific pattern
-- why the first entrant's timing is unrelated to demand, why the second entrant's timing responds to
proximity rather than population size, why the effect disappears after the second entrant, and why a
classical cross-sectional entry-threshold test need not show the same distance gradient the timing result
does (Section~\ref{sec:model}).

We are explicit throughout about what is well identified and what is not. The revenue result
(Section~\ref{sec:result1}) is robust across specifications and to excluding the largest outlier county. The
entry-timing result (Section~\ref{sec:result3}) is a real, placebo-tested pattern, but it is statistically
concentrated in three of the nine sample states, and a natural cross-sectional counterpart to the timing
story does not corroborate it. The model in Section~\ref{sec:model} is built to explain the timing pattern
including that last complication, not to paper over it.

\section{Data and Empirical Setting}
\label{sec:data}

The revenue results (Sections~\ref{sec:result1}--\ref{sec:result2}) use a county-year panel for Colorado
(2014--2025) and Washington (2015--2025), $n=706$ county-years with at least one open dispensary, built from
each state's cannabis-tax administrative data (Colorado Department of Revenue Socrata series; Washington
State Liquor and Cannabis Board fiscal-year sales reports) and a geocoded dispensary license registry used to
construct a current store count for each county-year. Montana's contrast case
(Section~\ref{sec:result2}) uses its Department of Revenue's monthly county sales reports, 2022--2025.
``Border'' counties are those within 70 miles of the nearest county in a state where recreational retail is
not yet legal, using each county's population-weighted centroid; we also report a continuous version of this
distance.

The entry-timing results (Section~\ref{sec:result3}) use a broader nine-state sample (California, Colorado,
Illinois, Massachusetts, Maine, New York, Ohio, Oregon, and Washington; 479 counties), each dispensary's own
opening date from a geocoded license registry, and each state's own retail-legalization date. This is a
different, larger sample than the two-state revenue panel, a distinction we return to in
Section~\ref{sec:result3} because it matters for how the timing result should be interpreted.

\section{Result 1: Revenue per Dispensary Is Higher in Border Counties}
\label{sec:result1}

Table~\ref{tab:a8_revenue_border} reports the core border-proximity specification. Counties within 70 miles
of a non-legal border earn significantly higher revenue per dispensary than interior counties of the same
two states: a coefficient of $+0.395$ (county-clustered $p=0.015$) with the full control set (log
population, year and state fixed effects), corresponding to roughly 48 percent higher revenue per store.
This is not driven by a single county -- excluding Sedgwick County, Colorado, a small, heavily touristed
county on the Nebraska border, the estimate falls only to $+0.302$ and remains significant. The continuous
version of the same relationship, log(distance to the nearest non-legal county), is negative as expected and
significant under heteroskedasticity-robust inference but not once standard errors are clustered by county;
we report both figures rather than the more favorable one alone.

\begin{table}[htbp]
\centering
\caption{Revenue per Dispensary: Border vs.\ Interior Counties (Colorado + Washington)}
\label{tab:a8_revenue_border}
\small
\begin{tabular}{p{5.6cm}cccc}
\toprule
Distance measure & Coef. & SE (HC1) & SE (cl.) & $p$ (cl.) \\
\midrule
Border $\le$70mi (binary), no controls & $+0.502$ & $0.101$ & $0.233$ & $0.031$ \\
Border $\le$70mi (binary), full controls$^{a}$ & $+0.395$ & $0.074$ & $0.163$ & $0.015$ \\
\quad excl.\ Sedgwick County, CO & $+0.302$ & $0.068$ & $0.147$ & $0.040$ \\
\midrule
log(distance), continuous, full controls$^{a}$ & $-0.182$ & $0.054$ & $0.118$ & $0.124$ \\
\bottomrule
\end{tabular}
\begin{minipage}{0.95\textwidth}
\vspace{0.3em}
\footnotesize
\textit{Notes:} Outcome is log(total county cannabis revenue $\div$ number of currently open dispensary
stores). County-year panel, Colorado and Washington, $n=706$ (71 counties, 10 border and 61 interior at the
70-mile threshold). $^{a}$Full controls: log(population), year fixed effects, state fixed effects. The
continuous-distance coefficient is significant under heteroskedasticity-robust (HC1) inference
($p<0.001$, not shown) but loses significance once standard errors are clustered by county; we report the
more conservative clustered $p$-value here and discuss this openly rather than lead with the more favorable
HC1 figure. The binary 70-mile result is robust to both inference methods and to excluding the single
largest outlier county.
\end{minipage}
\end{table}


Figure~\ref{fig:a8_levels} shows the raw difference in levels, and Figure~\ref{fig:a8_dose} shows the
continuous relationship directly as a residualized binned scatter: revenue per store falls smoothly as
distance to the border grows, not merely at the 70-mile cutoff.

\begin{figure}[htbp]
\centering
\includegraphics[width=0.62\textwidth]{figures/a8_fig1_revenue_levels.png}
\caption{Mean revenue per dispensary, border vs.\ interior counties, Colorado and Washington, 2014--2025.}
\label{fig:a8_levels}
\end{figure}

\begin{figure}[htbp]
\centering
\includegraphics[width=0.72\textwidth]{figures/a8_fig2_doseresponse.png}
\caption{Revenue per dispensary against distance to the nearest non-legal county, both residualized on
log(population), year fixed effects, and state fixed effects. Points are 10-bin (decile) means; point size
reflects the number of observations in each bin.}
\label{fig:a8_dose}
\end{figure}

\section{Result 2: A License-Cap Mechanism}
\label{sec:result2}

Why would extra demand show up as higher revenue per store rather than more stores? The natural candidate is
that the number of dispensaries is not free to adjust. Both Colorado and Washington constrain retail entry:
Washington capped the statewide number of licenses at 334 from the outset, allocating them by lottery in
oversubscribed jurisdictions \citep{WSLCB2014LotteryPR}; Colorado imposes no statewide cap but gives local
governments full authority under Amendment 64 to ban retail marijuana outright or to regulate its number,
and most Colorado municipalities have exercised this option to prohibit it entirely. \citet{HollenbeckUetake2021}
show directly that Washington's licensing restrictions produced substantial market power among incumbent
retailers. Under the classic entry-and-market-size logic of \citet{BresnahanReiss1990} and
\citet{BresnahanReiss1991}, if the number of firms a market supports is constrained below what demand would
otherwise justify, additional demand must appear as higher revenue per existing firm rather than more firms.

Montana offers a within-sample test of this logic, since it is the one state in our data with no statewide
numerical cap on dispensary licenses. If the license-cap story is right, Montana's border counties should
show the \emph{opposite} pattern: captive demand absorbed through more stores (higher revenue per capita)
rather than through higher revenue at each existing store.

\begin{table}[htbp]
\centering
\caption{The License-Cap Mechanism: Montana as an Uncapped Contrast Case}
\label{tab:a8_montana}
\small
\begin{tabular}{lcc}
\toprule
Outcome & Coefficient (border effect) & $p$-value (county-cl.) \\
\midrule
\multicolumn{3}{l}{\textit{Colorado + Washington (statewide/local license caps bind)}} \\
log(revenue per store) & $+0.395$ & $0.015$ \\
log(revenue per capita) & $+0.046$ & $0.833$ \\
\midrule
\multicolumn{3}{l}{\textit{Montana (no statewide numerical cap on dispensary licenses)}} \\
log(revenue per capita), 2022--2025 panel & $+0.654$ & $0.036$ \\
log(revenue per store), 2025 cross-section & $+0.012$ & $0.958^{a}$ \\
\bottomrule
\end{tabular}
\begin{minipage}{0.95\textwidth}
\vspace{0.3em}
\footnotesize
\textit{Notes:} $^{a}$HC1 standard error; the Montana revenue-per-store cross-section has too few border
counties ($n=23$) for county-clustered inference. Colorado and Washington cap the number of retail licenses
(Washington explicitly, statewide, at 334 from the outset; Colorado indirectly, through local governments'
authority to ban or cap retail marijuana under Amendment 64). Montana imposes no equivalent statewide cap.
The pattern reverses exactly as the license-cap mechanism predicts: where entry is capped, the border effect
appears in revenue per store and not per capita; where entry is free to expand, it appears in revenue per
capita (more stores absorbing the demand) and not per store.
\end{minipage}
\end{table}


This is exactly what we find. In Colorado and Washington, the border effect is significant in revenue per
store ($p=0.015$) and null in revenue per capita ($p=0.833$). In Montana, the pattern reverses: the border
effect is significant in revenue per capita ($p=0.036$) and a precise null in revenue per store ($p=0.958$).
A single coefficient could be a coincidence; a predicted reversal that shows up in exactly the state whose
institutional structure differs in exactly the relevant way is harder to dismiss as one.

\section{Result 3: The Second Entrant Arrives Faster Near the Border}
\label{sec:result3}

Revenue levels do not tell us whether sellers perceived the border opportunity and acted on it, or whether
demand simply sorted itself out after entry happened on its own schedule. We use each dispensary's own
opening date to ask about entry \emph{speed} rather than the eventual count, via Cox proportional-hazards
models stratified by state, over the nine-state entry-timing sample described in Section~\ref{sec:data}.

\begin{table}[htbp]
\centering
\caption{Entry Speed by Entrant Rank: Cox Proportional-Hazards Models}
\label{tab:a8_hazard_rank}
\small
\begin{tabular}{lccc}
\toprule
Entrant rank & Hazard ratio, log(distance) & $p$-value & Events / at-risk \\
\midrule
1st (state launch $\to$ 1st dispensary) & $0.97$ & $0.803$ & 346 / 479 \\
\textbf{2nd (1st $\to$ 2nd dispensary)} & $\mathbf{0.70}$ & $\mathbf{0.003}$ & 271 / 346 \\
3rd (2nd $\to$ 3rd dispensary) & $0.85$ & $0.266$ & 224 / 271 \\
4th (3rd $\to$ 4th dispensary) & $0.90$ & $0.496$ & 184 / 224 \\
\bottomrule
\end{tabular}
\begin{minipage}{0.95\textwidth}
\vspace{0.3em}
\footnotesize
\textit{Notes:} Cox models stratified by state, controlling log(population); nine-state entry-density
universe (CA, CO, IL, MA, ME, NY, OH, OR, WA; 479 counties). Hazard ratio $<1$ means faster entry closer to a
non-legal border. The effect is significant only for the second entrant.
\end{minipage}
\end{table}


The pattern is specific to the second entrant. There is no relationship between border distance and how
quickly a county gets its \emph{first} dispensary after its state's own retail launch (hazard ratio $0.97$,
$p=0.80$) -- consistent with regulatory and capital-raising delay, not demand, governing first entry. But
conditional on a county already having a first dispensary, a second one follows significantly faster the
closer that county is to a non-legal border (hazard ratio $0.70$, $p=0.003$): each additional log-mile of
distance lowers the hazard of a second entrant by roughly 30 percent. The effect weakens at the third entrant
and is gone by the fourth. Figure~\ref{fig:a8_hazard} plots all four ranks together.

\begin{figure}[htbp]
\centering
\includegraphics[width=0.68\textwidth]{figures/a8_fig3_hazard_by_rank.png}
\caption{Hazard ratio on log(distance to non-legal border), by entrant rank, with 95\% confidence
intervals. A ratio below 1 means faster entry closer to the border.}
\label{fig:a8_hazard}
\end{figure}

This pattern is consistent with -- though narrower than -- the general finding in the entry literature that a
market's second or third entrant matters disproportionately more than later ones
\citep{BresnahanReiss1990,BresnahanReiss1991}. Two further checks discipline what kind of effect this is.

\begin{table}[htbp]
\centering
\caption{Second-Entrant Hazard: What It Responds To}
\label{tab:a8_mechanism}
\begin{tabular}{lcc}
\toprule
Specification & Hazard ratio & $p$-value \\
\midrule
Distance to nonlegal border (main measure) & $0.70$ & $0.003$ \\
Captive nonlegal population within 70mi & $1.02$ & $0.111$ \\
\midrule
Placebo: distance to nearest county in \emph{any} state & $0.92$ & $0.478$ \\
\quad Nonlegal-border distance, raced against the placebo & $0.64$ & $0.003$ \\
\quad Placebo distance, raced against nonlegal-border distance & $1.15$ & $0.308$ \\
\bottomrule
\end{tabular}
\begin{minipage}{0.95\textwidth}
\vspace{0.3em}
\footnotesize
\textit{Notes:} All rows are the rank-2 (1st$\to$2nd entrant) specification from
Table~\ref{tab:a8_hazard_rank}. The effect responds to \emph{proximity} to a nonlegal jurisdiction, survives
a placebo test against generic (legal-status-blind) border proximity, but does not scale with the
\emph{size} of the population behind that border -- consistent with a travel-cost threshold rather than a
continuously demand-elastic response, formalized in Section~\ref{sec:model}.
\end{minipage}
\end{table}


The effect survives a placebo test: distance to the nearest county in \emph{any} other state, regardless of
legal status, is not significant on its own and loses a horse race against the nonlegal-border-specific
measure. This rules out a generic ``border towns compete faster'' story. But the effect does \emph{not} scale
with the \emph{size} of the population behind the border -- a captive-population measure is insignificant
alone and wrong-signed when raced against distance. Whatever this is, it responds to proximity, not to the
size of the market being revealed.

Two further limitations deserve equal billing with the finding itself. First, the statistical significance
of the pooled result is concentrated in Illinois, Massachusetts, and Ohio -- three states outside this
paper's core two-state revenue sample, with licensing regimes (merit-scored lotteries, locally negotiated
host-community agreements, and staged medical-to-adult-use conversion, respectively) different from Colorado
and Washington's. Colorado alone shows a same-direction but statistically insignificant version of the
pattern; Washington shows none. Second, a classical cross-sectional entry-threshold comparison -- the
population a county needs to support a second store, following \citet{BresnahanReiss1990} -- does not
corroborate the timing result: border counties need a \emph{larger}, not smaller, population increase to
reach a second store (a threshold ratio of $1.67$ versus $1.27$ in interior counties). We take this
seriously rather than set it aside, and the model below is built specifically to address it.

\section{Result 4: A Simple Entry Game for the Second Entrant}
\label{sec:model}

We now ask whether a simple, formal model can reproduce all three timing facts in Section~\ref{sec:result3}
at once: no effect at the first entrant, a significant effect at the second, no effect by the fourth, and no
corresponding shift in the cross-sectional population threshold. A model that reproduces only the headline
coefficient is not informative; one that also explains why the threshold test looks the way it does is a
real test of the story.

\subsection{Setup}

Consider a county $i$ with an unknown demand type $\theta_i \in \{L, H\}$. Type $H$ means residents of the
neighboring non-legal jurisdiction actually cross the border in meaningful numbers to buy legally; type $L$
means they do not, even if the county is geographically close, because of idiosyncratic factors such as road
access, local enforcement on the non-legal side, or consumer habits that a simple distance measure cannot
capture. Neither the analyst nor, initially, any potential entrant observes $\theta_i$ directly. The prior
probability of the high-demand type is
\begin{equation}
\Pr(\theta_i = H) = p(d_i), \qquad p'(d_i) < 0,
\end{equation}
decreasing in distance $d_i$ to the nearest non-legal county: proximity makes captive demand \emph{more
likely}, not \emph{certain} or \emph{proportionally larger}.

\begin{assumption}[Regulatory delay]
\label{as:delay}
Opening a dispensary requires clearing a licensing, zoning, and capital-raising process that takes random
time $\tau$, independent of $\theta_i$ and $d_i$.
\end{assumption}

\begin{assumption}[First-entry profitability]
\label{as:first}
Local, in-state demand alone is sufficient to make a first dispensary profitable in expectation, regardless
of $\theta_i$: $E_{\theta}[\pi_1(\theta_i, P_i)] > F$ for both types, where $P_i$ is local population and $F$
is the fixed cost of entry.
\end{assumption}

Assumption~\ref{as:first} is not imposed to fit the data; it is consistent with the model's own prediction
below and with a fact we already document in Section~\ref{sec:result1}: border counties support a first
dispensary at a smaller local population than interior counties (35,436 versus 44,602 residents in our data),
suggesting the demand bonus from a nearby border, even in expectation before it is confirmed, already lowers
the bar for a first entrant.

Under Assumptions~\ref{as:delay}--\ref{as:first}, the first entrant opens as soon as $\tau$ resolves,
regardless of $\theta_i$. This immediately delivers the first empirical fact: \textbf{first-entry timing does
not depend on $d_i$.}

Once open, the first dispensary's realized sales, foot traffic, and local reputation constitute a public
signal. We model this as a revealing event that arrives at rate $\mu_H$ if $\theta_i = H$ and never (or at a
negligible rate $\mu_L \approx 0$) if $\theta_i = L$ -- a type-$H$ market generates enough activity to be
locally noticeable within a random but bounded time; a type-$L$ market never does. A risk-neutral second
entrant from a competitive fringe monitors each county and enters as soon as the revealing signal arrives,
provided $\pi_2(H, P_i) > F > \pi_2(L, P_i)$: entering only once a high-demand market is confirmed, and never
absent that confirmation, is sequentially rational given a prior $p(d_i)$ too low, on its own, to justify
entry without the signal. This constitutes a perfect Bayesian equilibrium of the second-stage entry game: the
second entrant's threshold strategy (enter iff signal observed) is optimal given its beliefs, and its beliefs
(certainty of type $H$ upon the signal; unrevised prior absent it) are consistent with Bayes' rule on the
equilibrium path.

\subsection{Comparative statics}

\begin{proposition}
\label{prop:main}
Under Assumptions~\ref{as:delay}--\ref{as:first} and the equilibrium described above:
\begin{enumerate}
  \item First-entry timing does not depend on $d_i$.
  \item The hazard of second entry, averaged over the (econometrically unobserved) type $\theta_i$, is
        approximately $h_2(d_i) \approx p(d_i)\,\mu_H$, so that $\partial h_2/\partial d_i < 0$: second entry
        is faster the closer a county is to the border.
  \item If $p(d_i)$ falls off sharply beyond a travel-cost reach $\bar{d}$ rather than declining smoothly
        with the population behind the border, then conditional on being within reach, second-entry speed
        should respond weakly to the \emph{size} of that population.
  \item Once $\theta_i$ is revealed by the second entrant's arrival, no further private information remains
        to be resolved; entry beyond the second store follows an ordinary, $d_i$-independent process common
        to both types.
  \item A cross-sectional comparison of the population needed to support a second store, conditional on a
        county already having reached two stores, implicitly conditions on $\theta_i = H$ -- since only
        type-$H$ counties reach $N_i=2$ in this model. Conditional on $\theta_i=H$, the population threshold
        is a property of $\pi_2(H, \cdot)$ alone and need not vary systematically with $d_i$: $d_i$ affects
        the \emph{probability} of being selected into the high-demand group and the \emph{speed} of that
        selection being revealed, not the profit function of an already-confirmed high-demand market.
\end{enumerate}
\end{proposition}

Points 1 through 4 reproduce Table~\ref{tab:a8_hazard_rank} and Table~\ref{tab:a8_mechanism} directly: no
effect at rank 1, a significant effect at rank 2, insensitivity to population size conditional on reach, and
fade-out by rank 4. Point 5 is the model's account of the entry-threshold non-finding in
Section~\ref{sec:result3}: it is not evidence against the timing mechanism, because the timing mechanism
operates through \emph{selection and revelation speed}, not through a shift in the profit function of markets
that are already confirmed to be high-demand. We flag this as the model's own prediction, not a
post-hoc rescue of an inconvenient result -- the model would be false if second entry did not respond to
$d_i$ at all, and it does.

\begin{figure}[htbp]
\centering
\includegraphics[width=0.68\textwidth]{figures/a8_fig4_model_hazard.png}
\caption{Illustrative model-implied second-entry hazard, $p(d_i)\,\mu_H$, as a function of distance to the
border, using a logistic reach function for $p(d_i)$. The vertical line marks the 70-mile threshold used
throughout the empirical analysis.}
\label{fig:a8_model}
\end{figure}

The model is deliberately close to two existing traditions rather than a novel construction.
\citet{Jovanovic1982} models firms learning their own type through realized performance and exiting or
surviving accordingly; our second entrant instead learns \emph{a rival's} type through the rival's realized
performance, a small variation on the same Bayesian-learning logic. \citet{ChamleyGale1994} model investment
timing as a game of strategic delay driven by an informational externality -- players wait to see what others
reveal. Our second entrant is the limiting case in which the second mover simply reacts once the externality
resolves rather than strategically timing its own delay; a genuinely strategic extension, in which several
potential second entrants each wait for the others to reveal information, is a natural next step we do not
pursue here.

\section{Discussion}

We want to be precise about what this exercise does and does not establish. The revenue result
(Section~\ref{sec:result1}) is our most robust finding and the one we would lead with in any submission. The
license-cap mechanism (Section~\ref{sec:result2}) is supported by a real within-sample contrast, not merely
asserted, though it rests on one state's evidence and a modest sample. The entry-timing result
(Section~\ref{sec:result3}) is real and placebo-tested but statistically narrower than the revenue result --
concentrated in three states outside our core sample -- and we do not think it should be described as a
central contribution of this line of work. The model in Section~\ref{sec:model} is a reduced-form-motivated
theoretical account, not a structurally estimated one: we do not fit $\mu_H$, $p(d_i)$, or $F$ to the data,
and the hazard-ratio comparative statics it delivers are qualitative predictions that the empirical patterns
satisfy, not parameter estimates the model and data jointly identify. A fully structural version of this
model -- estimating $p(d_i)$ and $\mu_H$ by maximum likelihood from the hazard data directly -- is a natural
and, we think, tractable extension, but it is a larger undertaking than this paper attempts.

\section{Conclusion}

Cannabis retailers in U.S.\ counties near a non-legal state border earn more revenue per dispensary, and this
appears to operate through a license-cap channel that we can test directly using Montana's uncapped market as
a contrast. A second, narrower fact -- that a county's second dispensary opens faster the closer that county
is to the border, while its first does not -- is consistent with a simple entry game in which regulatory
delay governs first entry and Bayesian learning about a rival's revealed success governs the second. The
model's contribution is not to claim a new economic mechanism -- the general regularity that a market's
second entrant matters disproportionately is well established \citep{BresnahanReiss1990,BresnahanReiss1991}
-- but to show that a standard learning-based account of \emph{why} can reproduce every feature of our
timing evidence, including the one cross-sectional test that initially appeared to contradict it.

\newpage
\bibliographystyle{apalike}
\bibliography{references}

\end{document}
